% ##############################################################################
\section{Code Quality}

% Related documentation files:
% - helpers_root/docs/tools/linter/all.developing_linter.how_to_guide.md
%   How to develop custom linters
% - helpers_root/docs/tools/linter/all.amp_fix_md_links.explanation.md
%   Example linter for markdown links
% - helpers_root/docs/tools/all.precommit.how_to_guide.md
%   Pre-commit hooks setup and usage
% - helpers_root/docs/tools/all.code_coverage.how_to_guide.md
%   Code coverage measurement
% - helpers_root/docs/tools/all.codecov.how_to_guide.md
%   Codecov integration for coverage tracking
% - helpers_root/docs/tools/all.code_coverage_subprocess.how_to_guide.md
%   Coverage in subprocess testing
% - helpers_root/docs/tools/all.import_check.reference.md
%   Import cycle detection tools
% - helpers_root/docs/tools/git/all.git_hooks.how_to_guide.md
%   Git hooks for quality enforcement
% - helpers_root/docs/tools/git/all.git_hooks.reference.md
%   Complete git hooks reference
% - helpers_root/docs/tools/git/all.gitleaks.how_to_guide.md
%   Secret detection with gitleaks
% - helpers_root/docs/tools/git/all.gitleaks.reference.md
%   Gitleaks configuration reference
% - helpers_root/docs/code_guidelines/all.coding_style.how_to_guide.md
%   Coding style standards
% - helpers_root/docs/code_guidelines/all.coding_style_guidelines.reference.md
%   Complete style reference
% - helpers_root/docs/code_guidelines/all.code_review.how_to_guide.md
%   Code review process and quality gates
% - helpers_root/docs/code_guidelines/all.submit_code_for_review.how_to_guide.md
%   Preparing code for review
% - helpers_root/docs/code_guidelines/all.type_hints.how_to_guide.md
%   Type hints usage for better quality
% - helpers_root/docs/build_system/all.linter_gh_workflow.explanation.md
%   Automated linting in CI/CD
% - helpers_root/docs/build_system/all.semgrep_workflow.explanation.md
%   Static analysis with Semgrep

% ==============================================================================
\subsection{Motivation}
Code quality is of paramount importance for building maintainable software
systems. As codebases grow in size and complexity, maintaining consistent
standards, detecting architectural issues, and ensuring comprehensive test
coverage become increasingly challenging. Manual code reviews alone cannot
scale to catch all violations, architectural degradation, or quality
regressions before they enter the codebase. To address these challenges, we
developed an integrated code quality system that enforces standards
automatically, provides immediate feedback to developers, and maintains
visibility into code health through comprehensive measurement and reporting.

% ==============================================================================
\subsection{Summary}

% TODO(gp_ai3): The text should correspond to the topics here and convert to latex

Our integrated code quality system combines multiple enforcement layers to
maintain high standards throughout the development lifecycle.

- The pluggable linter framework provides both automatic fixes and quality
  gates, ensuring consistent formatting and catching common errors.

- Import cycle detection prevents architectural degradation by enforcing
  acyclic dependencies.

- Pre-commit hooks provide immediate feedback at commit time, catching issues
  before they enter the repository.

- Comprehensive coverage tracking with structured test suites and
  enforced thresholds ensures all code paths remain tested as the codebase
  evolves.

- This multi-layered
approach creates a robust quality enforcement system that scales with team
size and codebase complexity. By automating enforcement and providing
immediate feedback, we reduce the burden on code reviewers while maintaining
consistent standards across all contributions. The system is designed to
be transparent, providing clear diagnostics when checks fail and maintaining
audit trails of all validations. Together, these components enable teams to
move quickly while maintaining high standards.

% Make it easy to create packages
% Ensure code has same formatting (black, ruff, ...)
% Catch errors statically using linting
% Ensure code follows the same higher level rules (e.g., using AI)
% - interface style (input, output, inout)
% - testing
% Keep code in order of dependencies (interfaces first and then the rest in
% order to make reading it easy for a human)
% Docstring formatting
% Convention for branch naming to make it easy to refer to bug tracking context
% Documentation convention (diataxis)

% ==============================================================================
\subsection{The Linter Framework}
At the core of our code quality enforcement is a pluggable linting framework
that checks and modifies code to improve readability, detect bugs, and enforce
coding standards. Unlike monolithic linters, our system is built as a
collection of focused, composable linting scripts that can be executed
independently or as an integrated pipeline.

% ------------------------------------------------------------------------------
\subsubsection{Architecture}
Each linter is implemented as a Python class that inherits from
\texttt{linters.action.Action}
and implements a standardized interface:
\begin{itemize}
  \item \textbf{\texttt{\_execute()}}: Processes a file and returns a list
    of detected violations with descriptions

  \item \textbf{\texttt{check\_if\_possible()}}: Determines if the linter
    can run in the current environment

  \item \textbf{\texttt{run()}}: Orchestrates execution across multiple files
\end{itemize}

Linters are categorized into two types based on their behavior:
\begin{itemize}
  \item \textbf{Modifying actions}: Automatically fix issues by rewriting
    files (e.g., code formatters, import sorters, whitespace normalizers)

  \item \textbf{Non-modifying actions}: Report violations without changing
    files (e.g., static analysis tools, style checkers, complexity
    analyzers)
\end{itemize}

This separation enables different execution strategies: modifying linters
run during development to automatically maintain standards, while non-modifying
linters act as quality gates in CI/CD pipelines.

% ------------------------------------------------------------------------------
\subsubsection{Comprehensive Linting Rules}

Our linting framework integrates both third-party tools and custom rules
tailored to our coding conventions. The modifying linters include:

\begin{itemize}
  \item \textbf{Code formatting}: ty, Black for consistent Python formatting,
    Prettier for Markdown files

  \item \textbf{Import management}: isort for import ordering, autoflake
    for removing unused imports

  \item \textbf{Structural consistency}: Class method ordering,
    separating line normalization, whitespace standardization

  \item \textbf{Documentation}: Docstring formatting, table of contents
    generation for notebooks
\end{itemize}

Non-modifying linters enforce additional constraints:

\begin{itemize}
  \item \textbf{Static analysis}: flake8, pylint, and mypy for catching bugs
    and type errors

  \item \textbf{Policy enforcement}: File size limits, naming conventions,
    merge conflict detection

  \item \textbf{Convention checking}: TODO format validation, import
    pattern verification, shebang requirements
\end{itemize}

% ------------------------------------------------------------------------------
\subsubsection{Integration and Extensibility}
Adding a new linting rule requires creating a Python script in the
\texttt{linters/} directory, registering it in \texttt{linters/base.py}, and
adding unit tests. This modular design allows teams to customize enforcement
without modifying core infrastructure. Linters can be executed selectively
based on file types, directories, or specific development phases, providing
flexibility while maintaining consistency.

% ==============================================================================
\subsection{Import Cycle Detection}
Circular dependencies between modules represent a significant
architectural anti-pattern that degrades code maintainability, complicates
testing, and can lead to subtle runtime errors. As codebases grow organically,
developers may inadvertently introduce import cycles that violate
architectural boundaries. To maintain clean architecture, we developed
automated tools for detecting and visualizing circular dependencies.

% ------------------------------------------------------------------------------
\subsubsection{Dependency Analysis Tools}
Our system provides two complementary tools for import analysis:
\begin{itemize}
  \item \textbf{show\_imports}: A visualization tool that generates dependency
    graphs showing relationships between modules and packages. It
    supports multiple visualization modes including file-level
    dependencies, directory-level aggregation, external package
    dependencies, and cycle-only views. The tool can limit analysis
    depth, filter by scope, and export results in various image formats.

  \item \textbf{detect\_import\_cycles}: An enforcement tool that analyzes
    the dependency graph and fails if circular dependencies are detected.
    When cycles are found, it logs all participating files in each cycle,
    enabling developers to quickly identify and resolve architectural
    violations.
\end{itemize}

Both tools are built on Python's import machinery and require that all
directories containing Python files are proper modules (containing
\texttt{\_\_init\_\_.py}).
This constraint ensures comprehensive analysis and prevents silent
failures from malformed package structures.

% ------------------------------------------------------------------------------
\subsubsection{Integration into Development Workflow}
Import cycle detection runs as part of our continuous integration pipeline,
preventing pull requests with circular dependencies from being merged. Developers
can also run cycle detection locally before committing, receiving
immediate feedback on architectural violations. The visualization tools aid
in understanding complex dependency relationships and planning refactoring
efforts to break cycles. By enforcing acyclic dependencies automatically,
we maintain architectural integrity as the codebase evolves.

% ==============================================================================
\subsection{Package Structure and Import Conventions}

As Python codebases scale, the organization of modules and their import
relationships becomes critical to maintainability. Without clear conventions,
imports become fragile during refactoring, circular dependencies accumulate,
and understanding module relationships becomes increasingly difficult.
We developed a package organization system that prevents these problems
through clear conventions, hierarchical structure, and automated enforcement.

% ------------------------------------------------------------------------------
\subsubsection{Goals of Packages}

The goal of creating packages is to simplify imports from clients,
hide which file contains the actual code to enable reorganization without
changing client code, organize code in related units, and avoid import
loops by enforcing that there are no import loops in any module or
among modules. For example, importing from a package:

\begin{verbatim}
import dataflow.core as dtfcore
dtfcore.ArmaGenerator(...)
\end{verbatim}

is cleaner than importing the specific file:

\begin{verbatim}
import dataflow.system.source_nodes as dtfsysonod
dtfsysonod.ArmaGenerator(...)
\end{verbatim}

% ------------------------------------------------------------------------------
\subsubsection{Rules for Imports}

We follow rules to avoid import loops:

\begin{itemize}
  \item \textbf{Code inside a package imports files directly}: Within
    \texttt{datapull/common/data/client/}, code imports sibling files
    as \texttt{import datapull.common.data.client.abstract\_im\_clients
    as dcabimcl}, not the package itself. This prevents self-referential
    cycles.

  \item \textbf{Code from other packages imports the package}: External
    code imports \texttt{import datapull.common.data.client as dcdc}
    and uses \texttt{dcdc.AbstractImClient(...)}, not internal files.
    This hides implementation details and enables internal reorganization.

  \item \textbf{No statically detectable import loops}: We disallow
    any import loop that can be detected by inspecting code without
    executing it. This guarantees no dynamic import loops exist. The
    invoke command \texttt{lint\_detect\_cycles} checks for violations
    and runs automatically in continuous integration.

  \item \textbf{Only module-level imports}: Imports must occur at module
    level, not inside functions. We do not accept using local imports
    to break import loops, as they hide dependencies from static analysis.

  \item \textbf{Controlled package creation}: We avoid creating too
    many packages since organizing and maintaining them has overhead.
    We create packages only when the benefits justify the costs.

  \item \textbf{Canonical import aliases}: Each package specifies a
    short import alias in its \texttt{\_\_init\_\_.py} file, constructed
    from package path initials and kept to eight characters or fewer.
    For example, \texttt{datapull.common.data.client} becomes
    \texttt{dcdc}, \texttt{dataflow.core} becomes \texttt{dtfcore},
    and \texttt{market\_data} becomes \texttt{mdata}. The alias is
    documented in a docstring and used consistently. Linters verify
    compliance.
\end{itemize}

\paragraph{Unit Test Imports}
Unit tests can use either package imports when testing exported functionality
or direct file imports when testing internal implementation. Both are
acceptable since tests are both external clients and validators of
implementation details. However, test files are always leaf nodes and
never import from each other. Common test code belongs in library
files like \texttt{foobar\_example.py} or \texttt{test/foobar\_test\_case.py}.

% ------------------------------------------------------------------------------
\subsubsection{Package Hierarchy and Cycle Prevention}

To prevent import cycles, we enforce that certain packages do not depend
on other packages. This creates a directed acyclic graph at the package
level.

The hierarchy consists of distinct layers: \texttt{helpers} contains
low-level utilities depending only on external libraries, \texttt{core}
provides project-specific utilities depending only on \texttt{helpers},
\texttt{dataflow} implements the computational framework depending on
\texttt{core} and \texttt{helpers}, and application packages
(\texttt{oms}, \texttt{market\_data}, \texttt{im}) implement business
logic depending only on lower layers.

Within \texttt{helpers} itself, we maintain a micro-hierarchy among
utilities to prevent cycles in foundational code. The hierarchy progresses
from \texttt{hwarnings} and \texttt{hlogging} through \texttt{hdbg},
then \texttt{hintrospection} and \texttt{hprint}, to \texttt{henv},
\texttt{hsystem}, and \texttt{hio}, finally to \texttt{hgit} which
depends on system operations. A library can import only from preceding
or same-level libraries without creating cycles.

Attempts to introduce upward dependencies (such as \texttt{helpers}
importing \texttt{core}) are rejected during code review. This disciplined
approach maintains architectural clarity as the codebase evolves.

% ------------------------------------------------------------------------------
\subsubsection{Anatomy of a Package}

Every package directory contains \texttt{\_\_init\_\_.py} serving
multiple purposes: marking the directory as a Python package, documenting
the canonical import alias in a docstring, and optionally exporting
the public interface. A typical structure:

\begin{verbatim}
"""
Market data access and processing.

Import as:

import market_data as mdata
"""

from market_data.abstract_market_data import AbstractMarketData
__all__ = ["AbstractMarketData"]
\end{verbatim}

While \texttt{import *} is generally prohibited, it is acceptable within
\texttt{\_\_init\_\_.py} files for controlling the package's public
interface.

% ------------------------------------------------------------------------------
\subsubsection{Special Cases and Enforcement}

\paragraph{Unit Test Imports}
Unit tests can use either package imports when testing public interfaces
or direct file imports when testing internal implementation. Both
perspectives have merit since tests are both external clients and
validators of implementation details. However, test files must be leaf
nodes, they never import from each other. Shared test functionality
belongs in dedicated library files like \texttt{foobar\_example.py}
or \texttt{test/foobar\_test\_case.py}.

\paragraph{Enforcement Mechanisms}
Import conventions are enforced through multiple complementary mechanisms.
Static cycle detection tools analyze the dependency graph and flag
violations. Package structure verification ensures all directories
with Python files include \texttt{\_\_init\_\_.py}. Canonical alias
checking via linters verifies documented aliases are used. Automated
CI checks prevent pull requests with violations from merging. Code
review ensures new packages fit the hierarchy and architectural boundaries
are respected.

% ------------------------------------------------------------------------------
\subsubsection{Benefits}

This structured approach delivers concrete benefits. Internal file
reorganization within packages does not break external code, enabling
continuous improvement without cascading breakages. The package hierarchy
makes layering explicit, preventing ad-hoc coupling between components.
The combination of hierarchical structure and import rules makes circular
dependencies structurally difficult. Canonical aliases and consistent
patterns reduce cognitive load, letting developers focus on writing
code rather than debating import style. As the codebase grows, the
package structure provides natural boundaries for understanding subsystems
in isolation.

The primary trade-off is upfront structure. Creating new packages
requires thoughtful consideration of hierarchy placement, interface
design, and alias documentation. However, this investment pays dividends
as maintaining discipline from the start is far easier than untangling
dependencies later.

% ==============================================================================
\subsection{Pre-Commit Hook System}
Automated checks at commit time form the first line of defense against
quality regressions. While CI/CD pipelines catch issues before merging,
pre-commit hooks provide immediate feedback before committing, reducing
context switching and accelerating iteration cycles. Our pre-commit hook
system enforces essential invariants before allowing commits to succeed.

% ------------------------------------------------------------------------------
\subsubsection{Enforced Checks}
The pre-commit hook (\texttt{pre-commit.py}) executes a series of checks
automatically when developers attempt to commit changes:
\begin{itemize}
  \item \textbf{Branch verification}: Prevents accidental commits directly
    to protected branches (e.g., master)

  \item \textbf{Author validation}: Ensures Git user name and email are properly
    configured according to organizational standards

  \item \textbf{File size limits}: Blocks commits containing files exceeding
    size thresholds to prevent repository bloat

  \item \textbf{Python compilation}: Validates that all modified Python files
    compile successfully, catching syntax errors immediately

  \item \textbf{Secret detection}: Scans for accidentally committed secrets,
    credentials, or sensitive data using gitleaks integration

  \item \textbf{Forbidden patterns}: Detects merge conflict markers and other
    problematic patterns in committed files
\end{itemize}
If any check fails, the commit is aborted and developers receive a clear
diagnostic message indicating which check failed and why. Only when all checks
pass does the commit proceed.

% ------------------------------------------------------------------------------
\subsubsection{Commit Message Enhancement}
In addition to pre-commit validation, a commit-msg hook (\texttt{commit-msg.py})
enforces commit message conventions and automatically appends metadata
about executed checks. Each commit message includes a footer documenting
which pre-commit checks were run and passed, providing an audit trail of
quality enforcement. This transparency helps during code review and debugging
by clearly indicating what validations occurred.

% ------------------------------------------------------------------------------
\subsubsection{Installation and Configuration}
Git hooks are installed automatically when developers activate the thin
environment via \texttt{setenv.sh}, ensuring consistent enforcement
across all team members. The hooks can be manually installed or removed using
\texttt{install\_hooks.py} for explicit control. While hooks can be bypassed
using \texttt{git commit --no-verify} for exceptional circumstances, this
practice is discouraged and requires conscious decision-making.
Repository administrators can disable hooks for specific projects via
\texttt{repo\_config.yaml} when necessary, though the default is
universal enablement.

% ==============================================================================
\subsection{Code Coverage Tracking and Enforcement}
Maintaining comprehensive test coverage across a growing codebase requires
more than just writing tests, it demands visibility, automation, and
enforcement. Our integration with Codecov provides a system-wide view of
test coverage, structured into fast, slow, and superslow test suites. This
setup ensures that all code paths are exercised and that test coverage regressions
are identified early and reliably. The coverage system works in concert
with the linting and pre-commit systems to form a comprehensive quality
enforcement framework.

% ==============================================================================
\subsection{Structured Coverage by Test Category}

We categorize coverage tests into three suites based on runtime and
scope:

\begin{itemize}
  \item Fast tests run frequently (e.g., daily) and provide immediate feedback
    on high-priority code paths

  \item Slow tests cover broader logic and data scenarios

  \item Superslow tests are comprehensive, long-running regressions
    executed on a weekly cadence or on-demand
\end{itemize}

Each suite produces its own coverage report, which is flagged and
uploaded independently to Codecov, enabling targeted inspection and carryforward
of data when some suites are skipped.

% ==============================================================================
\subsection{CI Integration and Workflow Behavior}

Coverage reports are generated and uploaded automatically as part of our
CI pipelines. The workflow:

\begin{itemize}
  \item Fails immediately on critical setup errors (e.g., dependency or
    configuration issues)

  \item Continues gracefully if fast or slow tests fail mid-pipeline,
    but surfaces those failures in a final gating step

  \item Treats superslow failures as critical, immediately halting the
    workflow
\end{itemize}

This behavior ensures resilience while preventing silent test
degradation.

% ==============================================================================
\subsection{Enforced Thresholds and Quality Gates}

Coverage checks are enforced at both project and patch levels:

\begin{itemize}
  \item Project-level threshold: Pull requests fail if overall coverage
    drops beyond a configured margin (e.g., \textgreater1\%)

  \item Patch-level checks: Changes are required to maintain or improve
    coverage on modified lines

  \item Flags and branches: Checks are scoped per test suite and only
    enforced on critical branches
\end{itemize}

Together, these gates maintain coverage integrity while avoiding noise from
unrelated code paths.

% ==============================================================================
\subsection{Visibility and Developer Experience}

Codecov is integrated tightly into the developer workflow:

\begin{itemize}
  \item PRs show inline coverage status and file-level diffs

  \item Optional summary comments detail total coverage, changes, and affected
    files

  \item Reports can be viewed in Codecov's UI or served locally as HTML

  \item Carryforward settings retain historical data when full test suites
    aren't executed
\end{itemize}

Developers can also generate and inspect local reports for any test suite
using standard coverage commands.

% ==============================================================================
\subsection{Best Practices and Operational Consistency}

To ensure effective usage:

\begin{itemize}
  \item Coverage is always uploaded---even if tests fail---ensuring no blind
    spots

  \item Developers are encouraged to monitor coverage deltas in PRs

  \item The system defaults to global configuration, but supports fine-tuning
    via repo-specific overrides

  \item Weekly reviews of coverage trends and flags help spot regressions
    and low-tested areas
\end{itemize}

% ##############################################################################
\section{Unit test framework}

% Related documentation files:
% - helpers_root/docs/tools/unit_test/all.write_unit_tests.how_to_guide.md
%   Core guide to writing unit tests with hunitest framework
% - helpers_root/docs/tools/unit_test/all.run_unit_tests.how_to_guide.md
%   Running and organizing unit tests
% - helpers_root/docs/tools/all.code_coverage.how_to_guide.md
%   Measuring test coverage
% - helpers_root/docs/tools/all.codecov.how_to_guide.md
%   Coverage reporting and enforcement
% - helpers_root/docs/tools/notebooks/all.run_jupyter_notebook.how_to_guide.md
%   Testing Jupyter notebooks
% - helpers_root/docs/tools/notebooks/all.jupyter_notebook.how_to_guide.md
%   Working with notebooks in tests
% - helpers_root/docs/build_system/all.pytest_allure.explanation.md
%   Test reporting with Allure
% - helpers_root/docs/build_system/all.pytest_allure.how_to_guide.md
%   Setting up Allure for test visualization
% - docs/oms/broker/ck.generate_broker_test_data.how_to_guide.md
%   Example of generating test data for complex systems
% - docs/oms/broker/ck.replayed_ccxt_exchange.explanation.md
%   Mocking external dependencies for testing

% ==============================================================================
\subsection{Summary}
% TODO(gp_ai3): The text should correspond to the topics here and convert to latex

%Our unit testing framework combines principled philosophy with practical
%tooling.
%
%- By inheriting from \texttt{hunitest.TestCase}, tests gain
%access to enhanced assertions, golden file testing, and directory
%management utilities.
%
%- Categorizing tests by speed enables fast feedback cycles
%without sacrificing comprehensive regression testing.
%
%- The opinionated
%approach to mocking—mock only external dependencies—keeps tests focused
%on behavior and maintainable over time. 
%
%- Coverage measurement and enforcement
%ensure quality gates without creating busywork.
%
%- check_string vs assert_equal
%  - Support for fuzzy testing
%
%- Test organization and conventions
%
%- Mocking conventions
%
%- Special support for
%Jupyter notebooks extends testing to interactive, exploratory workflows.
%This testing infrastructure is foundational to our development system. It
%enables confident refactoring, documents expected behavior, and catches regressions
%early.
%
%- Test coverage
%
%- The consistency of patterns and conventions reduces cognitive
%load, allowing developers to focus on writing tests that matter rather than
%debating how to structure them.

% Different levels of testing
% Where to save tests

% ==============================================================================
\subsection{Testing Philosophy and Motivation}
Testing needs to be a first-class concern in software development.

We view unit tests as
executable specifications of behavior. Each test documents what the
system should do, provides a safety net for refactoring, and serves as
living documentation for developers. Good tests improve code design by forcing
clear interfaces and reducing coupling between components.

% ==============================================================================
\subsection{The hunitest.TestCase Framework}
At the core of our testing infrastructure is \texttt{hunitest.TestCase},
a custom test base class that extends Python's standard \texttt{unittest}
framework with specialized functionality for our development workflows.
Every test class in our codebase inherits from \texttt{hunitest.TestCase},
ensuring consistent testing patterns and access to enhanced testing
utilities. The framework provides several categories of functionality:
\begin{itemize}
  \item \textbf{Enhanced assertions}: Methods like \texttt{assert\_equal()}
    that provide superior diff output compared to standard \texttt{assertEqual()},
    especially for large strings and data structures

  \item \textbf{Golden file testing}: The \texttt{check\_string()}
    method for regression testing against frozen reference outputs
    stored on disk

  \item \textbf{Directory management}: Automatic creation and cleanup of
    test directories for inputs, outputs, and scratch space

  \item \textbf{Text processing}: Utilities for fuzzy matching,
    whitespace normalization, and memory address purification in test assertions

  \item \textbf{Notebook execution}: Specialized support for testing Jupyter
    notebooks as executable artifacts
\end{itemize}

% ------------------------------------------------------------------------------
\subsubsection{Basic Test Structure}
Every test method follows a consistent three-section structure that
makes tests self-documenting and easy to understand:

\begin{lstlisting}[language=Python]
def test_something(self) -> None:
  """
  Brief description of what this tests.
  """
  # Prepare inputs.
  input_data = create_test_data()
  # Run test.
  actual = function_under_test(input_data)
  # Check outputs.
  expected = "expected_result"
  self.assert_equal(str(actual), str(expected))
\end{lstlisting}

This structure enforces clarity: setup is separate from execution, which
is separate from verification. When a test fails, developers can immediately
identify which phase failed and why. The mandatory docstring and type
hints further improve readability and maintainability.

% ==============================================================================
\subsection{Golden File Testing with check\_string}
For complex outputs or large data structures, comparing against inline expected
values becomes impractical. Golden file testing solves this by storing
reference outputs in versioned files alongside the test code.
The \texttt{check\_string()} method implements this pattern.

% ------------------------------------------------------------------------------
\subsubsection{How check\_string Works}
When a test calls \texttt{self.check\_string(actual)}, the framework:
\begin{enumerate}
  \item Converts the actual output to a string representation

  \item Computes a path based on the test class and method name (e.g.,
    \texttt{test/outcomes/TestFoo1.test\_bar1/output/test.txt})

  \item Compares the actual output against the stored golden file

  \item Reports any differences using sophisticated diff algorithms

  \item Provides a command to update the golden file if the change is intentional
\end{enumerate}

This approach has several advantages over inline assertions:
\begin{itemize}
  \item Large outputs don't bloat the test code

  \item Diffs are easier to review in version control

  \item Tests remain readable even for complex outputs

  \item Updating expectations is a single command: \texttt{pytest --update\_outcomes}

  \item Historical reference outputs are preserved in git history
\end{itemize}

% ------------------------------------------------------------------------------
\subsubsection{When to Use check\_string vs assert\_equal}
The framework provides clear guidance:
\begin{itemize}
  \item Use \texttt{assert\_equal()} for: Simple comparisons (e.g.,
    \texttt{1 + 1 == 2}), small outputs (< 50 lines), values that should
    never change

  \item Use \texttt{check\_string()} for: Large outputs (> 50 lines),
    complex data structures, outputs that may evolve with the system,
    regression testing of behavior
\end{itemize}
Both methods support optional parameters for flexible comparison: \texttt{fuzzy\_match=True}
for whitespace-insensitive comparison, \texttt{purify\_text=True} for
stripping memory addresses and timestamps, and \texttt{dedent=True} for
automatic indentation normalization.

% ==============================================================================
\subsection{Test Organization and Conventions}
Consistent test organization reduces cognitive overhead and makes the codebase
easier to navigate. Our conventions cover naming, file placement, and
structural patterns.

% ------------------------------------------------------------------------------
\subsubsection{File and Directory Structure}
Tests live in \texttt{test/} directories adjacent to the code they test.
For a module \texttt{foo/bar/module.py}, the test file is \texttt{foo/bar/test/test\_module.py}.
Test artifacts (golden files, input data) are stored in subdirectories
following the pattern \texttt{test/TestClassName.test\_method\_name/}.
This co-location ensures tests are easy to find and maintain. When code moves,
tests move with it. When code is deleted, tests are deleted too.

% ------------------------------------------------------------------------------
\subsubsection{Naming Conventions}
Test class names clearly indicate what they test:
\begin{itemize}
  \item For a class \texttt{FooBar}: \texttt{TestFooBar1}

  \item For a function \texttt{process\_data()}: \texttt{Test\_process\_data1}

  \item For a method \texttt{FooBar.calculate()}: test method \texttt{TestFooBar1.test\_calculate1()}
\end{itemize}
Numbers (e.g., \texttt{TestFooBar1}, \texttt{test\_calculate1()}) are used
even for the first test to facilitate adding additional test classes or methods
later without renaming files and directories. Test methods use
descriptive names when there are few variations (\texttt{test\_calculate\_with\_empty\_input()})
or numbered names when testing multiple similar scenarios (\texttt{test\_calculate1()},
\texttt{test\_calculate2()}, \texttt{test\_calculate3()}).

% ------------------------------------------------------------------------------
\subsubsection{Helper Methods and DRY Principle}
When multiple test methods follow the same pattern with different inputs,
helper methods eliminate repetition:

\begin{lstlisting}[language=Python]
class TestFooBar1(hunitest.TestCase):
  def _helper(self, input_val, expected):
      # Run test.
      actual = foo_bar(input_val)
      # Check outputs.
      self.assert_equal(str(actual), str(expected))

  def test1(self) -> None:
      self._helper(input_val=1, expected="result1")

  def test2(self) -> None:
      self._helper(input_val=2, expected="result2")
\end{lstlisting}

This keeps tests DRY (Don't Repeat Yourself) while maintaining clarity about
what distinguishes each test case.

% ==============================================================================
\subsection{Test Categorization and Execution}
Tests are categorized by execution time to enable fast feedback cycles. The
framework defines three categories using pytest markers:

\begin{table}[h]
\centering
\begin{tabular}{|l|p{5cm}|p{5cm}|}
\hline
\textbf{Test Category} & \textbf{Examples of Tests} & \textbf{When to Run} \\
\hline
\textbf{Fast tests} (< 5 sec per test class) & Core functionality, business logic, unit-level tests & Run before every commit and in CI on every pull request \\
\hline
\textbf{Slow tests} (< 20 sec) & More complex scenarios, integration-level tests, tests involving external APIs (mocked) & Run daily and before releases \\
\hline
\textbf{Superslow tests} (< 5 min) & End-to-end workflows, production simulations, comprehensive regression tests & Run weekly or on-demand \\
\hline
\end{tabular}
\caption{Test categorization by execution time}
\label{tab:test_categories}
\end{table}

% ------------------------------------------------------------------------------
\subsubsection{Running Tests with invoke}
The invoke task system provides convenient commands for running each
category:

\begin{lstlisting}[language=bash]
# Run Fast Tests Only
invoke run_fast_tests
# Run Slow Tests
invoke run_slow_tests
# Run Superslow Tests
invoke run_superslow_tests
# Run with Coverage
invoke run_fast_tests --coverage
# Run Specific Test
invoke run_fast_tests --pytest-opts="path/to/test.py::TestClass::test_method"
\end{lstlisting}
Tests always execute inside Docker containers to ensure environmental consistency.
The framework automatically handles pytest options, timeout configuration,
and result reporting.

% ------------------------------------------------------------------------------
\subsubsection{Timeout and Retry Behavior}
Tests are subject to timeouts based on their category to prevent runaway
tests from blocking CI pipelines. The \texttt{pytest-timeout} plugin enforces
these limits. Additionally, \texttt{pytest-rerunfailures} automatically
retries flaky tests: fast tests are retried twice, slow and superslow
tests once. This balances reliability with execution time.

% ==============================================================================
\subsection{Mocking Philosophy and Practice}
Our approach to mocking is pragmatic and opinionated: mock only external
dependencies, never mock internal code. This philosophy stems from our end-to-end
testing preference and our focus on behavioral correctness.

% ------------------------------------------------------------------------------
\subsubsection{Mock Only External Dependencies}
We mock interactions with systems outside our control:
\begin{itemize}
  \item Third-party APIs (e.g., CCXT cryptocurrency exchange library)

  \item Cloud infrastructure (S3, AWS Secrets Manager, databases)

  \item External services (GitHub API, OpenAI, data providers)
\end{itemize}
For example, when testing code that fetches data from a cryptocurrency
exchange via CCXT, we mock the CCXT library, not our wrapper around it:

\begin{lstlisting}[language=Python]
@umock.patch.object(module.ccxt, "binance")
def test_fetch_ohlcv(self, mock_binance):
  mock_exchange = umock.Mock()
  mock_binance.return_value = mock_exchange
  mock_exchange.fetch_ohlcv.return_value = test_data
  # Now test our code's behavior
\end{lstlisting}
This ensures we're actually testing our code's logic, not the mock
itself.

% ------------------------------------------------------------------------------
\subsubsection{Do Not Mock Internal Code}
We never mock code within our own repositories. Doing so creates
maintenance burdens and tests implementation details rather than
behavior. If internal code is difficult to test without mocking, that's
a signal to refactor for better testability. This approach means tests exercise
real code paths, catching integration issues that mocking would hide. It
also makes refactoring safer: as long as the public interface remains
stable, internal restructuring doesn't require updating tests.

% ------------------------------------------------------------------------------
\subsubsection{Shared Mock Setup}
For test classes that require common mocks, we use \texttt{set\_up\_test()}
and \texttt{tear\_down\_test()} with pytest fixtures:

\begin{lstlisting}[language=Python]
class TestExchangeClient1(hunitest.TestCase):
  # Create patches as class variables
  ccxt_patch = umock.patch.object(module, "ccxt")
  secret_patch = umock.patch.object(module.hsecret, "get_secret")
  @pytest.fixture(autouse=True)
  def setup_teardown_test(self):
      self.set_up_test()
      yield
      self.tear_down_test()
  def set_up_test(self) -> None:
      self.ccxt_mock = self.ccxt_patch.start()
      self.secret_mock = self.secret_patch.start()
      self.secret_mock.return_value = {"key": "test"}
  def tear_down_test(self) -> None:
      self.ccxt_patch.stop()
      self.secret_patch.stop()
\end{lstlisting}

This pattern ensures mocks are consistently configured across all test
methods and properly cleaned up after each test.

% ==============================================================================
\subsection{Test Coverage Measurement and Enforcement}
Code coverage is measured using the \texttt{coverage.py} library
integrated with pytest via \texttt{pytest-cov}. Coverage is tracked
separately for fast, slow, and superslow test suites, providing visibility
into which code paths are exercised by which test categories.

% ------------------------------------------------------------------------------
\subsubsection{Running Coverage Analysis}
Coverage analysis is initiated with the \texttt{--coverage} flag:

\begin{lstlisting}[language=bash]
# Generate Coverage for Fast Tests
invoke run_fast_tests --coverage
# Generate Coverage for Specific Module
invoke run_fast_tests --coverage --pytest-opts="oms/"
# Combine Coverage From Multiple Runs
docker> coverage combine .coverage_fast .coverage_slow
docker> coverage report --include="oms/*" --omit="*/test_*.py"
\end{lstlisting}

The framework generates both terminal reports and browsable HTML
coverage reports (\texttt{htmlcov/index.html}) showing line-by-line coverage
with color-coded indicators for covered, uncovered, and partially covered
lines.

% ------------------------------------------------------------------------------
\subsubsection{Coverage Integration with CI/CD}
Coverage reports are uploaded to Codecov as part of our CI/CD pipelines.
Each test suite (fast, slow, superslow) produces a separate coverage flag,
enabling targeted analysis. Pull requests are gated on coverage thresholds:
project-level coverage must not decrease by more than 1\%, and modified
lines must maintain or improve coverage. This prevents coverage regression
while allowing flexibility for exploratory development.

% ------------------------------------------------------------------------------
\subsubsection{Interpreting Coverage Metrics}
We focus on behavioral coverage, not line coverage. The goal is not 100\%
line coverage but comprehensive testing of observable behaviors and critical
code paths. Some lines, like defensive assertions and error handling for
impossible states, may remain uncovered by design. Coverage reports guide,
rather than dictate, testing efforts. Low coverage on a module signals
potential risk, but the appropriate response depends on context:
critical trading logic demands exhaustive testing, while simple utility functions
may not.

% ==============================================================================
\subsection{Testing Jupyter Notebooks}
Jupyter notebooks present unique testing challenges. They combine code,
narrative, and visualizations in a format designed for interactive exploration,
not automated testing. Our framework treats notebooks as executable artifacts
that must be validated for correctness and reproducibility.

% ------------------------------------------------------------------------------
\subsubsection{Notebook Execution Framework}
Notebooks are tested by executing them in a controlled environment using
\texttt{run\_notebook.py}, a wrapper around \texttt{papermill} that injects
parameters, captures output, and validates results:

\begin{lstlisting}[language=bash]
# Parameterized Notebook Execution
/app/dev_scripts/notebooks/run_notebook.py \
--notebook path/to/notebook.ipynb \
--config_builder "config_builder_function" \
--dst_dir /tmp/results
\end{lstlisting}

This approach ensures notebooks execute without errors, produce expected
outputs, and remain compatible with current code versions.

% ------------------------------------------------------------------------------
\subsubsection{Notebook Test Organization}
Notebook tests follow the same naming conventions as regular tests. For
a notebook \texttt{analysis.ipynb}, the test is \texttt{test\_analysis\_notebook.py}.
The test class typically has a single test method that executes the
notebook and validates key outputs or side effects. Tests use golden file
testing to detect unintended changes in notebook output. When notebooks
generate plots, tables, or metrics, these are captured and compared against
reference versions stored in test outcomes.

% ------------------------------------------------------------------------------
\subsubsection{Debugging Failing Notebook Tests}
When a notebook test fails, the debugging workflow is:
\begin{enumerate}
  \item Run the failing test with \texttt{-s --dbg} to get detailed logs

  \item Extract the \texttt{run\_notebook.py} command from the logs

  \item Re-run with \texttt{--no\_suppress\_output} to see full notebook
    output

  \item Identify the failing cell and root cause from the detailed error
    report
\end{enumerate}

This workflow provides full visibility into notebook execution while
keeping test output concise during normal runs.

% ##############################################################################
\section{invoke workflows}

% Related documentation files:
% - helpers_root/docs/tools/all.invoke_workflows.how_to_guide.md
%   Core guide to using invoke for task automation
% - helpers_root/docs/tools/all.invoke_git_branch_copy.how_to_guide.md
%   Specific invoke task example for git operations
% - helpers_root/docs/tools/all.invoke_bash_print_tree.how_to_guide.md
%   Invoke task for directory visualization
% - docs/infra/ck.kaizen_infra_script.how_to_guide.md
%   Infrastructure invoke tasks
% - docs/infra/ck.kaizen_infra_python_script.how_to_guide.md
%   Python-based invoke workflows
% - docs/infra/ck.kaizen_infra_python_script.reference.md
%   Reference for infra invoke tasks
% - helpers_root/docs/onboarding/all.development.how_to_guide.md
%   Common invoke commands for daily development

% Invoke Task Automation: Centralized task registry using Python Invoke for common development operations (testing, linting, deployment).

% TODO(ai_gp2): Add Motivation

% TODO(ai_gp2): Add Summary

% ##############################################################################
\section{Code Repo Automation}

% Related documentation files:
% - helpers_root/docs/tools/dev_system/all.synchronize_gh_issue_labels.how_to_guide.md
%   Automated label synchronization
% - helpers_root/docs/tools/dev_system/all.synchronize_gh_settings.how_to_guide.md
%   Repository settings synchronization
% - helpers_root/docs/tools/all.sync_gh_projects.how_to_guide.md
%   GitHub Projects synchronization
% - helpers_root/docs/tools/git/all.invite_gh_contributors.explanation.md
%   Automating contributor management
% - helpers_root/docs/tools/git/all.invite_gh_contributors.how_to_guide.md
%   How to manage contributors programmatically
% - helpers_root/docs/work_organization/all.github_projects_process.reference.md
%   GitHub Projects structure and automation
% - helpers_root/docs/work_organization/all.pr_workflow.explanation.md
%   Pull request automation workflow
% - helpers_root/docs/work_organization/all.issue_workflow.explanation.md
%   Issue lifecycle automation

% ==============================================================================
\subsection{Motivation}

As organizations grow, so does the surface area of their code repo
ecosystem---more repositories, more contributors, and more operational
complexity. Without strong conventions and enforcement, labels diverge
across projects, and team boards lose structural consistency. This leads
to disjointed workflows, unclear ownership, and fragmented planning. To mitigate
this, we introduced a declarative automation system for GitHub metadata.
It enables centralized definitions for labels and project structures and
propagates them across the organization in a reproducible and scalable way.

We settle on git as version control system and GitHub as ...

Multiple alternative solutions are available that would 

% ==============================================================================
\subsection{Summary of Solutions}

- Make it easy to create composable Git modules for reuse and access control
- Ensure that each Git repo follows the same standards (e.g., unified labels,
  configuring GitHub projects)
- A declarative approach to repository settings

% ==============================================================================
\subsection{Standardized Label Infrastructure}

Issue labels form the foundation for triage, workflows, and reporting. At
scale, inconsistent naming, coloring, or descriptions can fragment
automation and reduce clarity. To address this, we maintain a
centralized manifest that defines the entire organization's label
taxonomy. Each label includes a name, color code, and description.
Repositories synchronize against this manifest using a containerized
process that:

\begin{itemize}
  \item Creates missing labels using manifest definitions

  \item Updates outdated labels (e.g., if the description or color
    changes)

  \item Backs up existing labels before applying changes

  \item Supports dry-run execution for visibility and confidence

  \item Optionally prunes unused labels not defined in the manifest
\end{itemize}

This guarantees all repositories speak a consistent language for issues and
pull requests, improving developer experience and enabling reliable
automation.

% ==============================================================================
\subsection{Template-Based Project Synchronization}

GitHub Projects (Beta) are a powerful tool for planning and tracking,
but manually configuring project fields and views across teams is
tedious and error-prone. To streamline setup and reduce drift, we
introduced a project templating system. Project metadata, such as fields
and views is defined in a canonical source project and synced into
destination projects. The system:

\begin{itemize}
  \item Adds missing global fields to ensure functional parity across projects

  \item Logs discrepancies in view structure (e.g., missing ``Backlog''
    or ``Current sprint'' views)

  \item Preserves existing configurations to avoid accidental data loss

  \item Operates in a dry-run mode to preview proposed changes safely
\end{itemize}

Due to current API limitations, view creation, layout configuration, and
field ordering cannot yet be automated. However, warnings are logged for
manual follow-up, and the system is designed to evolve as GitHub expands
its GraphQL support.

% ==============================================================================
\subsection{Design Principles and Workflow Behavior}

% TODO(ai_gp2): Merge these two things

These automation tools follow several core principles:

\begin{itemize}
  \item Declarative Configuration: Metadata is defined in structured
    YAML files stored in version control, serving as the source of truth

  \item Reproducible Execution: All sync operations run in isolated
    Docker containers, ensuring consistent behavior across machines and CI
    pipelines

  \item Non-Destructive by Default: No changes are applied unless
    explicitly confirmed, and current states can be backed up for rollback

  \item Extensibility: The architecture is modular and designed to
    incorporate future GitHub API enhancements, including view creation and
    layout syncing
\end{itemize}

The workflows support both interactive execution and automated pipelines,
allowing integration into CI/CD systems or local tooling as needed.

% ==============================================================================
\subsection{Declarative Repository Settings Synchronization}

As code base footprint grows, so has the need to manage repository
configuration such as metadata fields, merge policies, and branch protection
rules in a reliable and consistent manner across dozens of projects. Manual
configuration inevitably leads to drift, accidental misconfiguration, and
wasted effort during onboarding or project setup. To address this, we
introduced a \textit{declarative} synchronization system for repository
settings, inspired by the principles of infrastructure-as-code. Repository
policies—including key settings (description, default branch, visibility, merge
strategies, etc.) and fine-grained branch protection rules are defined
centrally in a YAML manifest. This manifest serves as the single source of
truth for desired state. A containerized executable can \textbf{export} the
current settings of any repository for inspection, backup, or review, and can
\textbf{apply} (sync) the manifest to one or more repositories, ensuring
alignment with organizational standards. The system also supports a dry-run
mode for safe preview of all proposed changes.

This approach delivers several benefits:
\begin{itemize}
  \item \textbf{Consistency:} All repositories adhere to the same
    baseline policies, reducing configuration drift and error.

  \item \textbf{Auditability:} Changes to settings are reviewed and
    version-controlled like code, supporting transparent change
    management.

  \item \textbf{Scalability:} New projects can be onboarded with a
    single command, inheriting standardized policies from day one.

  \item \textbf{Recoverability:} Before each sync, the system backs up the
    existing configuration, enabling rollbacks if needed.

  \item \textbf{Confidence:} Dry-run and logging capabilities provide
    full visibility into what changes would be made, supporting both
    local experimentation and safe automation in CI/CD pipelines.
\end{itemize}

% ##############################################################################
\section{CI/CD Automation}

% Related documentation files:
% - helpers_root/docs/work_organization/all.buildmeister.how_to_guide.md
%   Core Buildmeister role and responsibilities
% - helpers_root/docs/work_organization/all.inframeister.how_to_guide.md
%   Infrastructure-focused CI/CD monitoring
% - helpers_root/docs/work_organization/all.datapull_DAGmeister.how_to_guide.md
%   Data pipeline CI/CD monitoring
% - helpers_root/docs/build_system/all.pytest_allure.explanation.md
%   Allure test reporting in CI/CD
% - helpers_root/docs/build_system/all.pytest_allure.how_to_guide.md
%   Setting up Allure reports
% - helpers_root/docs/build_system/all.gitleaks_workflow.explanation.md
%   Gitleaks secret detection in CI/CD
% - helpers_root/docs/build_system/all.linter_gh_workflow.explanation.md
%   Linter automation in GitHub Actions
% - helpers_root/docs/build_system/all.semgrep_workflow.explanation.md
%   Semgrep security scanning workflow
% - docs/build/ck.github_actions.explanation.md
%   GitHub Actions CI/CD setup
% - docs/build/ck.gh_workflows.explanation.md
%   Workflow definitions and structure
% - docs/build/ck.gh_workflows.reference.md
%   Complete workflow reference
% - docs/build/ck.setup_github_actions.how_to_guide.md
%   Initial GitHub Actions setup
% - docs/build/ck.debug_github_actions.how_to_guide.md
%   Debugging CI/CD failures
% - docs/build/ck.github_runners.explanation.md
%   GitHub runners configuration
% - docs/build/ck.setup_slack_notification_for_actions.how_to_guide.md
%   Build notification setup

% ==============================================================================
\subsection{Motivation}

Automated test pipelines are essential, but without accountability, they
often fall into disrepair. The Buildmeister routine introduces a
rotating, human-in-the-loop system designed to enforce green builds,
identify root causes, and ensure high-quality CI/CD hygiene. This
mechanism aligns technical execution with team responsibility, fostering
a culture of operational ownership and co-operation.

% ==============================================================================
\subsection{Why It Matters}

The Buildmeister is a system of shared accountability. It transforms test
stability from an abstract ideal into a daily operational habit, backed by
clear expectations, defined processes, and human enforcement. By combining
automation with ownership, the team achieves sustainable reliability in a
complex, multi-repo ecosystem.

% ==============================================================================
\subsection{Core Responsibilities}

The Buildmeister is a rotating role assigned to a team member every one or two weeks. Their primary duties are:

\begin{itemize}
  \item Monitor build health daily via the Buildmeister Dashboard

  \item Investigate failures and ensure corresponding GitHub Issues are filed
    promptly

  \item Push responsible team members to fix or revert breaking code

  \item Maintain test quality by analyzing trends in reports

  \item Document breakage through a structured post-mortem log
\end{itemize}

The Buildmeister ensures the policy of: ``Fix it or revert within one hour.''

% ==============================================================================
\subsection{Handover and Daily Reporting}

The routine begins each day with a status email to the team detailing:

\begin{itemize}
  \item Overall build status (green/red)

  \item Failing test names and owners

  \item GitHub issue references

  \item Expected resolution timelines

  \item A screenshot of the Buildmeister dashboard
\end{itemize}

At the end of each rotation, the outgoing Buildmeister must confirm handover
by receiving a reply from the incoming one, ensuring continuity and awareness.

% ==============================================================================
\subsection{Workflow in Practice}

When a build breaks:

\begin{itemize}
  \item The team is alerted via Slack (\#build-notifications) through
    our GitHub Actions bot

  \item The Buildmeister triages the issue:

    \begin{itemize}
      \item Quickly reruns or replicates the failed tests if uncertain

      \item Blames commits to identify the responsible party

      \item Notifies the team and files a structured GitHub Issue
    \end{itemize}

  \item All information including test names, logs, responsible engineer
    are transparently shared and tracked
\end{itemize}

If the issue is not resolved within one hour, the Buildmeister must escalate
and disable the test with explicit owner consent or revert the offending change.

% ==============================================================================
\subsection{Tools and Analysis}

% ------------------------------------------------------------------------------
\subsubsection{Buildmeister Dashboard}

A centralized UI provides a real-time view of all builds across repos and
branches. It is the Buildmeister's daily launchpad.

% ------------------------------------------------------------------------------
\subsubsection{Allure Reports}

\begin{itemize}
  \item Every week, the Buildmeister reviews trends in skipped/failing
    tests, duration anomalies, and retry spikes

  \item This process:

    \begin{itemize}
      \item Surfaces hidden test instability

      \item Provides historical context to new breaks

      \item Enables preventive action before regressions cascade
    \end{itemize}
\end{itemize}

% ------------------------------------------------------------------------------
\subsubsection{Post-Mortem Log}

Every build break is logged in a shared spreadsheet, capturing:

\begin{itemize}
  \item Repo and test type

  \item Link to the failing GitHub run

  \item Root cause

  \item Owner and fix timeline

  \item Whether the issue was fixed or test was disabled
\end{itemize}

This living record forms the basis for failure mode analysis and future automation
improvements.
